\documentclass[a4paper, 12pt]{article}
\usepackage[portuges]{babel}
\usepackage[utf8]{inputenc}
\usepackage{amsmath}
\usepackage{indentfirst}
\usepackage{graphicx}
\usepackage[colorinlistoftodos]{todonotes}

\setlength{\marginparwidth}{2cm}
\begin{document}
    %\maketitle
    \begin{titlepage}
        \vspace{10pt}
        \begin{figure}[!ht]
            \centering
                \begin{center}
                    \includegraphics[width=3.5cm]{uninta.png}
                \end{center}
        \end{figure}
        \vspace{5pt}
    	\begin{center}
    		{\normalsize{CENTRO UNIVERSITÁRIO INTA - UNINTA}}
            \vspace{5pt}
        	\begin{center}
        		\normalsize{ATIVIDADE DISCURSIVA}
                \vspace{10pt}
                \vspace{71pt}
        		\large{Curso: Licenciatura em Pedagogia\\
                		Disciplina: Novas Tecnologas na Educação}
        		\vspace{140pt}
        	\end{center}
       	
        	\begin{flushleft}
        		\begin{center}
        			Aluno: Nome e Sobrenome
        	    \end{center}
        	\end{flushleft}
        	
        	\begin{center}
        	\vspace{\fill}
    		    \normalsize{Sobral, 24 de outubro de 2019}
        	\end{center}
    	\end{center}
    \end{titlepage}
%%%%%%%%%%%%%%%%%%%%%%%%%%%%%%%%%%%%%%%%%%%%%%%%%%%%%%%%%%%
\newpage
\tableofcontents
\thispagestyle{empty}

\newpage
\pagenumbering{arabic}

%%%%%%%%%%%%%%%%%%%%%%%%%%%%%%%%%%%%%%%%%%%%%%%%%%%%%%%%%
%%%%%%%%%%%%%%%%%%%%%%%%%%%%%%%%%%%%%%%%%%%%%%%%%%%%
\section{Resumo}

O objetivo deste experimento é a realização 

\section{Introdução}

O objetivo deste experimento é a realização 
\section{Desenvolvimento}

Para o experimento em questão foram utilizados dois geradores síncronos em que um atuava como motor para configurar um conjunto motor síncrono-gerador síncrono. A função do motor síncrono é a de fornecer potência mecânica ao eixo do gerador para que, então, seja possível realizar os ensaios no gerador síncrono. As placas de dados do motor e do gerador síncrono estão ilustrados nas Figuras 


\section{Considerações Finais}


análise dados
\section{Conclusão}

conclusao



\newpage
\section{Referências}

[1] Chapman, S.J. -- Electric Machinery Fundamentals, 4th Edition;

[2] Fitzgerald, A. E. -- Máquinas Elétricas, 2da Edição;

\end{document}